\documentclass[9.5pt]{article}
\usepackage[margin=0.6in]{geometry}
\usepackage{times,xcolor,titlesec,booktabs,hyperref}
\definecolor{maroon}{HTML}{500000}
\titleformat{\section}{\normalsize\bfseries\color{maroon}}{}{0pt}{}
\titlespacing{\section}{0pt}{5pt}{2pt}
\setlength{\parindent}{0pt}\setlength{\parskip}{3pt}
\hypersetup{colorlinks=true,linkcolor=maroon,urlcolor=maroon}
\pagestyle{empty}
\begin{document}
{\color{maroon}\rule{\linewidth}{3pt}}
{\bfseries\large\color{maroon} Data checkpoint: Aggie Course Scout}\hfill Carlos Luna-Pena, CSCE 470
{\color{maroon}\rule{\linewidth}{1pt}}

\section*{Where the data came from}
I pulled the 2026-2027 Texas A\&M undergraduate and graduate catalogs from \url{https://catalog.tamu.edu/}. I read \texttt{robots.txt} first and stayed out of disallowed paths. I crawled with four workers, a 0.2 second delay per worker, and three retries, under the user agent \texttt{AggieCourseScout/1.0 (CSCE 470 class project; clmoon2@tamu.edu)}. I keep every raw page as gzipped HTML.

\section*{Counts and hashes}
\begin{tabular}{@{}lr@{}}
\toprule
Department index pages requested & 383\\
Pages with HTTP 200 & 383\\
Pages with zero course blocks & 8\\
HTTP 404 or empty & 0\\
Unique level + code records & 10,082\\
Undergraduate / graduate & 5,228 / 4,854\\
Departments represented & 264\\
\bottomrule
\end{tabular}

\smallskip
Output file: \texttt{data/processed/courses.jsonl}, 8.0 MB.\\
SHA-256: \texttt{a2cfbcece529332fce6cb4117953be9bbd250385843fec6a0650d0291c4f8ca2}.\\
Manifest: \texttt{data/processed/crawl\_manifest.json}, 383 rows with URL, status, byte count, and page SHA-256.

\section*{What's in a record}
One JSON object per line. Fields: \texttt{id} (level plus code), \texttt{level}, \texttt{code}, \texttt{number}, \texttt{title}, \texttt{credits}, \texttt{description}, \texttt{prerequisites}, \texttt{corequisites}, \texttt{department}, \texttt{dept\_code}, \texttt{college}, \texttt{course\_level}, \texttt{url}, \texttt{source\_title}.

\section*{What I cleaned}
I stripped HTML with BeautifulSoup and normalized non-breaking and zero-width spaces. I collapsed whitespace, removed credit and meeting-hour boilerplate, split prerequisites out with a label regex, dropped cross-listed title prefixes, and deduped on level plus code. College is mapped by department prefix, because the catalog doesn't print college on a department page.

\section*{Samples (truncated)}
\texttt{\{"code": "CSCE 310", "title": "Database Systems", "credits": "3", "description": "Introduction to database systems; data models; ...", "prerequisites": "Grade of C or better in CSCE 221..."\}}\\
\texttt{\{"code": "MATH 151", "title": "Engineering Mathematics I", "credits": "4", "description": "Rectangular coordinates, vectors, functions ...", "prerequisites": "..." \}}

\section*{What this can't do}
This is a snapshot. Courses added after 2026-11-22 are absent, and removed courses remain. The catalog has no meeting times, seats, instructors, or modalities, so the tool can't answer ``when does this meet.'' College is hand-mapped and 536 records in professional or unplaced prefixes use ``Other/Professional.'' Honestly, a dozen or so prerequisite strings with odd punctuation come out cut off. I left them in rather than hand-fix them. Research, seminar, and zero-credit records are kept because they are real catalog entries, which pushes the count above the proposed 6,000 to 9,000 range.

\end{document}
